\chapter{Confirmatory results: the pre-registered verdicts}
\label{ch:results}

\begin{quote}\itshape
This chapter in brief: both families were tuned under the frozen 50-trial budget, the test
fleet was evaluated exactly once, and the five pre-registered hypotheses were adjudicated.
Three confirmed, two refuted --- and the two surprises (detection flipping to the AI after
tuning; physics features refusing to help) teach more than the expected wins. Every number
here is confirmatory; every earlier number in chapters 6--7 was exploratory and is superseded
for decision-making purposes.
\end{quote}

\keyidea{The split verdict \emph{is} the message: learning wins timing and prognosis, physics wins nothing it was promised, and the confusable-isolation wall binds both families equally because it is informational, not algorithmic. Every refutation here is a result, not a defect.}


\section{Protocol executed as frozen}
\label{sec:results-protocol}

Per \code{docs/prereg-v1.md}: 50 Optuna trials per family
(\code{scripts/tune_f5.py}, TPE seed 0, objectives on validation only, per-task winners from
the frozen pool), then one --- and only one --- evaluation of the tuned configurations on the
20 test engines (\code{scripts/f5\_confirm.py}), then the verdict machinery of \cref{sec:stats-math} (one-sided exact
McNemar for paired binary outcomes, Wilcoxon signed-rank paired by engine, Holm correction
across the tested hypotheses, BCa intervals and Cliff's-delta effect sizes). Two execution notes, both disclosed in
the scripts and the decision register: no traditional trial met the FA\,$\leq$\,1 validation
budget, so its detection winner was selected by the lexicographic infeasibility rule (fewest
false alarms, then highest recall); and AI models fit on training engines only, which keeps
the conformal calibration on validation statistically valid. The full F5 execution ---
100 tuning trials plus the confirmatory pass --- cost about 36 minutes on the Apple M5
reference machine (measured at execution; the resumable campaign design prevents a meaningful
cold-start re-benchmark).

\begin{figure}[htbp]
  \centering
  \includegraphics[width=0.62\textwidth]{optuna_convergence}
  \caption{The symmetric tuning budget in action: best validation RUL RMSE against Optuna
  trial, both families given the same 50 trials. Each search converges within its budget; the
  comparison is between two tuned practitioners, not a tuned one and a default. Auto-generated.}
  \label{fig:optuna}
\end{figure}

\begin{table}[htbp]
  \centering\small
  \caption{The five pre-registered verdicts (single confirmatory pass, test split,
  Holm-corrected). Auto-generated from \texttt{data/processed/f5/verdicts.json}.}
  \label{tab:verdicts}
  % AUTOGENERATED by scripts/make_report_assets.py — do NOT edit by hand
\begin{tabular}{>{\raggedright\arraybackslash}p{0.16\textwidth}>{\raggedright\arraybackslash}p{0.15\textwidth}>{\raggedright\arraybackslash}p{0.41\textwidth}r}
\toprule
Hypothesis & Verdict & Confirmatory evidence (AI vs traditional) & $p$ (Holm) \\
\midrule
  H1 (detection) & CONFIRMED & recall 0.48 vs 0.13; delay 499 vs 6033 cy; FA 1=1 & 0.008 \\
  H2 (confusable isolation) & REFUTED & 0.31 vs 0.31 (n=13; discordant 4/4) & 0.64 \\
  H3 (RUL prognosis) & CONFIRMED & RMSE 1694/858 vs 3816/1981 cy (50/90\,\% life); $\delta$=0.89 & 8.6e-06 \\
  H4 (physics hybrid) & REFUTED & hybrid worse at 10/25/100\,\% (2852 vs 1772 cy at 10\,\%) & --- \\
  H5 (uncertainty) & CONFIRMED & coverage 0.883 vs 0.983 (nominal 0.90); half-width 1957 vs 11509 cy & --- \\
\bottomrule
\end{tabular}

\end{table}

\section{H1 confirmed: the story is the tuning, not the model}
\label{sec:res-h1}

The v0 chapters called the AI's detection failure ``a design gap, not a threshold gap''
(\cref{sec:det-sweep}) --- and the frozen search space contained exactly the suspected lever:
the detector's feature-window lengths. Given 50 trials, the tuner stretched the averaging
windows and the whitened-residual detector went from 0.06 recall (exploratory) to
\textbf{0.48 confirmatory recall at a 499-cycle median delay} --- twelve times earlier than
the tuned traditional detector, at the same false-alarm count (one each), McNemar Holm-$p$
0.008. Meanwhile the traditional side's tuned selection --- picked by the frozen
infeasibility rule on a validation split with only two clean engines --- generalized poorly
(0.13 test recall, worse than its untuned v0). Two transferable lessons: window length is
the first knob any change detector lives or dies by; and threshold selection on small clean
validation samples is fragile for \emph{any} method --- budget clean engines accordingly.

\begin{figure}[htbp]
  \centering
  \includegraphics[width=0.5\textwidth]{detection_delays}
  \caption{Detection lead time: one dot per detected episode, with the median bar. The AI
  detects far more episodes (11 vs 3) and far sooner (median 499 vs 6\,033 cycles after
  onset) at equal false alarms. Auto-generated.}
  \label{fig:det-delays}
\end{figure}

\section{H2 refuted --- the observability wall binds everyone}
\label{sec:res-h2}

On the thirteen confusable test episodes the tuned families tie exactly: 31\,\% each, and
the disagreement pattern is perfectly symmetric (four episodes only the AI gets right, four
only the traditional rule does). The wall the influence-matrix geometry predicted in
\cref{sec:icm} --- fault signatures 1.3$^{\circ}$ apart in a rank-3 measurement space ---
binds both families equally. The transferable conclusion is bought hardware, not software:
\textbf{if confusable-fault isolation matters operationally, add gas-path instrumentation}
(the extended sensor set multiplies the minimum signature angle by five,
\cref{tab:icm-grid}); no amount of modeling recovers information the sensors never captured.

\begin{figure}[htbp]
  \centering
  \includegraphics[width=0.92\textwidth]{isolation_confusion}
  \caption{Isolation confusion matrices on the test episodes: true fault (rows) versus
  predicted (columns), traditional (left) and AI (right). Neither concentrates on the
  diagonal --- both scatter faults onto confusable neighbours (e.g.\ toward \texttt{fan.eta}
  and \texttt{hpc.eta}), the mechanism behind the H2 tie. Auto-generated.}
  \label{fig:iso-confusion}
\end{figure}

\section{H3 and H5 confirmed --- prognosis is where the AI pays}
\label{sec:res-h3h5}

The tuned GRU confirms the exploratory pattern with pre-registered force: RMSE
1\,694\,/\,1\,042\,/\,858 cycles at 50/70/90\,\% of life against the tuned margin
extrapolation's 3\,816\,/\,4\,544\,/\,1\,981; NASA scores one to two orders of magnitude
lower; Cliff's $\delta = 0.89$ (a nearly clean sweep at engine level); Holm-$p$
$8.6\times10^{-6}$. And the uncertainty statement that matters to a planner --- H5 --- comes
with it: the conformal 90\,\% interval covers 88.3\,\% empirically at a $\pm1\,957$-cycle
half-width, versus the slope-band's over-covering 98.3\,\% at $\pm11\,509$. In operational
terms: \textbf{a maintenance window six times narrower, with honest stated confidence} ---
the difference between a plannable shop-visit forecast and a shrug.

\begin{figure}[htbp]
  \centering
  \includegraphics[width=0.74\textwidth]{rul_distribution}
  \caption{Per-engine confirmatory RUL error, traditional versus AI, at three life fractions
  (positive = optimistic; boxes span the interquartile range). The evidence is in the shape,
  not just the RMSE: the traditional errors are wide and skewed optimistic, the AI's are tight
  and near-unbiased. Auto-generated from the tuned test-fleet predictions.}
  \label{fig:rul-dist}
\end{figure}

\section{H4 refuted --- physics features are not free wins}
\label{sec:res-h4}

The stacked hybrid lost at every data fraction (2\,852 vs 1\,902 cycles at 10\,\%), exactly
as the exploratory runs showed and now on the pre-registered record. The mechanism bears
repeating for practitioners tempted by ``physics-informed'' as a label: smeared state
estimates are noise-bearing, mutually correlated features whose information the learner
already had. Physics injection needs an information channel the raw data lacks --- the two
untested mechanisms (twin-residual features, physics-constrained losses) remain open
questions, but the popular default of concatenating estimator outputs onto a network input
is now a documented negative on this benchmark.

\section{What an engineering organization should take from this}
\label{sec:res-industry}

Strip the statistics away and the five verdicts become five decisions a fleet manager can
act on. In order of how much money each moves:

\begin{enumerate}
  \item \textbf{Deploy learning for prognosis first.} The confirmatory RUL advantage ---
    $2.3$--$4.4\times$ depending on life fraction, with calibrated intervals --- is the clear
    business case (maintenance slotting, lease planning). State it against the right baseline:
    that factor is against the linear extrapolation an operator fields \emph{today}; against
    the best \emph{advanced} classical prognostic the margin narrows to $\sim$1.3$\times$
    (\cref{sec:f8-lrul}), which is still a win but a different sales pitch.
  \item \textbf{Detection: modern change-detection with tuned windows} --- the AI's win came
    from window search, a cheap, auditable mechanism, not from depth.
  \item \textbf{Isolation: buy sensors, not models.} The confusable wall is informational.
  \item \textbf{Distrust untuned comparisons in either direction} --- both families moved
    dramatically under the same budget; v0 scoreboards inverted.
  \item \textbf{Demand pre-registration of vendor claims.} This chapter's split verdict is
    what honest evaluation looks like; a clean sweep should raise suspicion.
\end{enumerate}

\section{Sim-to-real: the ranking survives on the canonical external benchmark}
\label{sec:sim-to-real}

Contribution C7 asks whether the method ranking is an artifact of our own simulator. The
check ran on NASA's \textbf{C-MAPSS FD001} --- the field's reference turbofan
run-to-failure benchmark (100 unseen test units) --- with a disclosed scope substitution:
the plan named N-CMAPSS, whose multi-gigabyte distribution violates the desktop-
reproducibility norm; FD001 answers the same ranking question at 5\,MB, and only the
prognosis task transfers (FD001 carries no event labels). Same method shapes as the main
study: a classical health-index linear extrapolation versus the same GRU family
(\code{scripts/sim_to_real.py}, three seeds).

Result: \textbf{RMSE 14.9 cycles (AI) versus 42.1 (traditional)} at the community-standard
cap of 130 --- the same direction and comparable magnitude of advantage as H3 on SynCFM56,
and an absolute GRU figure competitive with published FD001 results, on a benchmark this
project had no hand in generating. The prognosis conclusion is not an artifact of SynCFM56.
With this, the F5 gate closes in full.
